\documentclass[11pt]{ctexart}

% 页面与字体
\usepackage[a4paper,margin=2.4cm]{geometry}
\setCJKmainfont{SimSun}[BoldFont=SimHei]
\setCJKsansfont{SimHei}
\setCJKmonofont{KaiTi}

% 数学与表格
\usepackage{amsmath,amssymb}
\usepackage{booktabs}
\usepackage{array}
\usepackage{tabularx}
\usepackage{graphicx}
\usepackage[colorlinks=true,linkcolor=black,urlcolor=blue]{hyperref}

\newcommand{\dd}{\,\mathrm{d}}
\newcommand{\pd}[2]{\frac{\partial #1}{\partial #2}}

\title{\bfseries 问题 4 计算结果报告\\[0.3em]\large 坐标变换法（Landau 变换）求解移动边界问题}
\author{}
\date{}

\begin{document}
\maketitle
\vspace{-2.5em}

\section{任务与计算口径}

问题 4 要求：\textbf{考虑药材烘干过程中的尺寸变化（收缩），重新确定烘干所需时间（单位：h）}；按表 6 格式给出每隔 6 h、到中心距离每隔 0.5 cm 的水分浓度，以及每隔 60 s、距离每隔 0.1 cm 的完整结果 \texttt{result4.xlsx}。

问题 1--3 均假设半径 $R=2$ cm 固定不变。问题 4 引入附件 2 给出的\textbf{收缩半径经验数据} $R(t)$（2.0 cm $\to$ 1.198 cm），使扩散区域 $[0,R(t)]$ 随时间移动，成为\textbf{移动边界问题}。直接沿用固定网格会因边界运动而失效，故本文采用\textbf{坐标变换法（Landau 变换）}把移动域映射为固定域后离散求解。

交付口径：

\begin{center}
\begin{tabularx}{\textwidth}{lX}
\toprule
\textbf{项} & \textbf{取值} \\
\midrule
移动边界 & 附件 2 的 $R(t)$，中心差分（\texttt{np.gradient}）得 $R'(t)$，线性插值 \\
归一化坐标 & $E = r/R(t)\in[0,1]$，固定域有限体积网格 $N = 800$ \\
时间推进 & Crank--Nicolson（$\theta = 0.5$）隐式 + Picard/Newton 线性化，$\Delta t = 60$ s；初始 1 h 内用 2 s 子步消除启动振荡（见 3.5 节） \\
物性模型 & 附录 4 变参数模型（$\rho$、$c_p$、$k$、$D$ 均随 $C,T$ 变化） \\
对流项 & 坐标变换严格导出的 $+(E R'/R)\partial/\partial E$（中心差分，表面单侧） \\
终止判据 & $\max_r C < 0.15$ kg/kg（等价于中心 $C < 0.15$，见 V5） \\
\textbf{交付结果} & \textbf{$t_\text{dry} = 188\,610$ s = 52.39 h；分钟向上取整 3144 min = 52.40 h} \\
\bottomrule
\end{tabularx}
\end{center}

\section{控制方程与 Landau 变换}

\subsection{移动边界问题}

无限长圆柱、一维轴对称、热质弱耦合（推导同问题 1/2），但区域随 $R(t)$ 变化：
\begin{equation}
\rho c_p\pd{T}{t}=\frac{1}{r}\pd{}{r}\Big(k\,r\pd{T}{r}\Big),\qquad
\pd{C}{t}=\frac{1}{r}\pd{}{r}\Big(D\,r\pd{C}{r}\Big),\qquad r\in[0,R(t)]
\end{equation}

定解条件（$T_\infty,C_\infty$ 取附件 1，覆盖 0--14\,400 s 后维持末值）：
\begin{equation}
r=0:\ \pd{T}{r}=\pd{C}{r}=0;\qquad
r=R(t):\ -k\pd{T}{r}=h(T-T_\infty),\quad -D\pd{C}{r}=h_m(C-C_\infty)
\end{equation}

初值 $T_0=28\ ^\circ\mathrm{C}$、$C_0=2.55$ kg/kg。$h=25$ W/(m$^2\cdot$K)、$h_m=8\times10^{-7}$ m/s。

附录 4 的变参数经验式（问题 4 与问题 3 不同，物性随含水量、温度的依赖更强）：
\begin{equation}
\begin{aligned}
&\rho=760+90C,\quad c_p=1850+\frac{2150C}{C+1},\quad k=0.12+\frac{0.20C}{C+1}\\[2pt]
&D=4.2\times10^{-4}\,e^{-0.30/C}e^{-3850/T}\quad(T\ \text{以 K 计})
\end{aligned}
\end{equation}

\subsection{Landau 变换与对流项符号}

引入归一化坐标 $E=r/R(t)$，把移动边界 $[0,R(t)]$ 映射到\textbf{固定域} $[0,1]$。由链式法则：
\begin{equation}
E=\frac{r}{R(t)}\ \Rightarrow\ \Bigl.\pd{E}{t}\Bigr|_r=-\frac{rR'}{R^2}=-\frac{ER'}{R}
\quad\Longrightarrow\quad
\Bigl.\pd{C}{t}\Bigr|_r=\Bigl.\pd{C}{\tau}\Bigr|_E-\frac{ER'}{R}\,\pd{C}{E}
\end{equation}

代入物理方程，注意到 $r=ER$、$\partial/\partial r=(1/R)\partial/\partial E$，得变换后的传质方程（热方程完全同构）：
\begin{equation}
\pd{C}{\tau}=\frac{1}{R^2}\frac{1}{E}\pd{}{E}\Big(E\,D\,\pd{C}{E}\Big)\;+\;\frac{E R'}{R}\,\pd{C}{E}
\end{equation}

其中对流项 $\boxed{\;+(E R'/R)\,\partial C/\partial E\;}$ 是\textbf{坐标变换（移动网格）带来的}。

\begin{quote}
\textbf{关于符号的说明（重要）。} 题目给出的变换式写作 $-(E R'/R)\partial C/\partial E$，这是\textbf{符号笔误}。严格链式法则给出 $+E R'/R$：物理导数 $\partial C/\partial t|_r$ 与变换导数 $\partial C/\partial\tau|_E$ 之间相差 $-ER'/R$，移项到右端后对流项为正号。物理上，$R'<0$（收缩）时对流速度为负、把水分\textbf{向中心输送}，对应``收缩使水分浓缩、从而减缓浓度下降''，与干燥被推迟的直觉一致；若取负号则收缩反而加速干燥，物理错误。V3 的符号敏感性检验定量确认了这一点（见 6.3 节）。
\end{quote}

边界条件在 $E$ 域变为：中心 $E=0$ 轴对称，表面 $E=1$ 处 $-(k/R)\partial T/\partial E=h(T-T_\infty)$、$-(D/R)\partial C/\partial E=h_m(C-C_\infty)$。

\section{数值方法：四项创新处理}

\subsection{精确处理坐标变换的对流项}

对流项 $v_E\partial/\partial E$（$v_E=E R'/R$）\textbf{对温度、浓度两方程一致}地加入，内部用中心差分（二阶精度）、表面 $E=1$ 用单侧差分、中心 $E=0$ 系数恰为 0（$v_E(0)=0$）。不把对流项当作扩散的扰动，而是作为独立的守恒输运算子与扩散算子相加，因此变换后的格式在 $R'=0$ 时\textbf{逐位退化}为问题 1/2 的固定域格式（V6）。该对流项不可忽略：V3 显示它使 $t_\text{dry}$ 由 50.83 h（无对流）变为 52.40 h，占比约 3\%。

\subsection{$R(t)$ 由附件 2 数值微分得到}

附件 2 给出 145 个时刻的半径（$t=0$--259\,200 s，$R$ 2.0$\to$1.198 cm）。$R'(t)$ 用\textbf{中心差分}（\texttt{np.gradient}，端点单侧）求得，再用线性插值成连续函数 $R(t),R'(t)$。关键特征：\textbf{收缩高度前重}（图 16a）——前 6 h 半径由 2.0 掉到 1.374 cm，前 24 h 基本收缩到 1.204 cm，之后稳定在 $\sim$1.2 cm。因此对流项在 $t\lesssim24$ h 内活跃，与``表面快速失水驱动收缩''的物理一致（图 16b 的 $R'$ 在前段为负且量值最大，末段趋于 0）。

\subsection{网格随边界等比例缩放，保证空间分辨率一致}

$E$ 域取\textbf{均匀}有限体积网格，物理网格 $r_j(t)=E_j R(t)$ 随边界同步缩放。好处有二：

\begin{enumerate}
\item \textbf{相对分辨率全程一致}——不像固定网格那样在收缩末期出现``边界层被网格跨过''的问题；
\item \textbf{守恒权重恒为常数}——$V_j=(r_{j+1/2}^2-r_{j-1/2}^2)/(2R^2)=(E_{j+1/2}^2-E_{j-1/2}^2)/2$ 与 $t$ 无关，且 $\sum_jV_j=(1^2-0^2)/2=1/2$ \textbf{恒成立}，离散质量守恒是恒等式而非近似（V7）。
\end{enumerate}

\subsection{附录 4 变参数物性 + 移动边界 = 完整模型}

在移动域上按 Crank--Nicolson（$\theta = 0.5$）隐式推进。热方程对 $T$ 线性（系数冻结于 $C$），直接 Thomas 追赶求解；传质方程因 $D(C,T)$ 非线性，用 Picard/Newton 线性化（显式给出 $\partial D/\partial C = D\cdot0.30/C^2$ 的敏度项，构造雅可比）后追赶求解，每个时间步做 4 次 Picard 迭代。热质弱耦合交替：先解热场再用新温度更新 $D$ 求解浓度场。

\subsection{初始瞬变的启动子步（消除 Crank--Nicolson 启动振荡）}

表面干燥边界层在\textbf{数十秒内}形成（表面浓度在首分钟内下降 $\sim$0.4 kg/kg），这是相对 $\Delta t = 60$ s 而言的刚性瞬变。Crank--Nicolson（$\theta=0.5$）不是 L 稳定的，用大步长起步会在此产生非物理振荡：实测 $\Delta t = 60$ s 起步时，表面浓度在 $t=60$ s 处偏低 0.122 kg/kg（1.990 vs 收敛值 2.112）、$t=120$ s 又反超 0.065 kg/kg。振荡在 1 h 内衰减（$t=1800$ s 误差 $3\times10^{-4}$、3600 s 误差 $2\times10^{-6}$），且\textbf{只影响表面}——中心在最初 $\sim$2 h 内恒为 2.55，故 $t_\text{dry}$（由中心控制）完全不受影响。

处理办法：$t<3600$ s 的每个输出步内部用 2 s 子步（\texttt{march} 细分）推进，之后恢复正常 $\Delta t = 60$ s。启动子步几乎不增加计算量（1800 个子步），却把表面列的前几分钟误差由 $\sim$0.12 压到 $10^{-5}$ 量级，恢复单调下降。这与问题 3 报告``初始瞬变需要细分、后段放大''的处理一脉相承。

\section{产物清单}

\begin{center}
\begin{tabularx}{\textwidth}{lX}
\toprule
\textbf{文件} & \textbf{内容} \\
\midrule
\texttt{results/result4.xlsx} & 交付文件：单张 \texttt{Sheet1}，表头「时间\symbol{92}到药材中心的距离」+ 0, 0.1, \dots, 1.9 cm + 「药材表面」；自 60 s 起每 60 s 一行，共 3144 行 $\times$ 22 列，水分浓度保留 4 位小数 \\
\texttt{results/problem4\_tables.txt} & 表 6、烘干时长、终止时刻剖面 \\
\texttt{results/problem4\_verification.txt} & V1--V7 验证套件原始输出 \\
\texttt{results/p4\_fields.npz} & 时空场缓存（供绘图复用） \\
\texttt{figures/fig16--19} & 半径收缩、移动边界剖面、特征位置时程、验证对比 \\
\texttt{code/p4\_model.py} & Landau 变换移动边界守恒算子 + 求解器 + 二分终止 \\
\texttt{code/p4\_run.py} & 交付口径求解与结果输出 \\
\texttt{code/p4\_verify.py} & V1--V7 验证套件 \\
\texttt{code/p4\_figures.py} & 四张配图 \\
\bottomrule
\end{tabularx}
\end{center}

\section{计算结果}

\subsection{表 6　药材烘干过程的水分浓度（kg/kg，干基）}

\begin{center}
\begin{tabular}{cccccc}
\toprule
时间/h & 0 cm & 0.5 cm & 1.0 cm & 1.5 cm & 药材表面 \\
\midrule
6 & 1.9524 & 1.7828 & 1.2571 & —— & 0.5684 \\
12 & 0.8705 & 0.7739 & 0.4848 & —— & 0.2145 \\
18 & 0.4572 & 0.4112 & 0.2641 & —— & 0.1016 \\
24 & 0.3058 & 0.2793 & 0.1895 & —— & 0.0710 \\
30 & 0.2370 & 0.2189 & 0.1549 & —— & 0.0608 \\
36 & 0.1992 & 0.1854 & 0.1352 & —— & 0.0565 \\
42 & 0.1754 & 0.1641 & 0.1223 & —— & 0.0543 \\
48 & 0.1591 & 0.1494 & 0.1132 & —— & 0.0531 \\
\textbf{52.39} & \textbf{0.1500} & \textbf{0.1412} & \textbf{0.1081} & —— & \textbf{0.0525} \\
\bottomrule
\end{tabular}
\end{center}

注：\textbf{``——''表示该固定距离处材料已因收缩离开}（$r\ge R(t)$，该处无药材内部水分）。由附件 2，$R(6\text{h})=1.374$ cm，故 1.5 cm 处在 6 h 时即已因收缩而``脱离''药材——这是移动边界与固定网格（问题 1--3）在表格呈现上最直观的差别。

\subsection{烘干时长}

\begin{center}
\begin{tabular}{llll}
\toprule
口径 & $t$ / s & $t$ / h & 分钟口径 \\
\midrule
连续口径（二分到 $\sim$1 s） & 188\,610 & \textbf{52.3917} & 3143.5 min \\
\textbf{交付口径（分钟向上取整）} & 188\,640 & \textbf{52.4000} & \textbf{3144 min = 52 h 24 min} \\
\bottomrule
\end{tabular}
\end{center}

空间收敛（V1，$N=400$ 与 800 均为 52.4000 h）、时间收敛（V2，$\Delta t = 60/120$ s 均为 52.4000 h）均已达标，故取 \textbf{$t_\text{dry} \approx 52.4$ h}。

\subsection{终止时刻的径向剖面（$t = 188\,640$ s，$R = 1.2000$ cm）}

\begin{center}
\begin{tabular}{cccccccc}
\toprule
$r$ / cm & 0 & 0.1 & 0.2 & 0.3 & 0.4 & 0.5 & 0.6 \\
\midrule
$C$ & 0.1500 & 0.1496 & 0.1486 & 0.1469 & 0.1445 & 0.1412 & 0.1372 \\
\bottomrule
\end{tabular}
\end{center}

\begin{center}
\begin{tabular}{ccccccc}
\toprule
$r$ / cm & 0.7 & 0.8 & 0.9 & 1.0 & 1.1 & 表面(1.2) \\
\midrule
$C$ & 0.1321 & 0.1259 & 0.1181 & 0.1081 & 0.0936 & 0.0525 \\
\bottomrule
\end{tabular}
\end{center}

剖面单调递减，中心恰好落在判据上，表面 0.0525 已逼近环境平衡值 $C_\infty=0.0499$。

\subsection{结果的物理解读}

\begin{enumerate}
\item \textbf{收缩使烘干时长几乎减半。} 同一套附录 4 物性下，冻结 $R=2$ cm 的 $t_\text{dry}=129.1$ h，而计入收缩后为 52.4 h，缩短 \textbf{2.46 倍}（图 18b、图 19c）。机制有二：扩散路径由 2 cm 缩短到 $\sim$1.2 cm；扩散项被 $1/R^2$ 放大，$R$ 缩小 1.67 倍即放大 2.79 倍。这是移动边界模型与固定域模型之间\textbf{不可忽略的系统性差别}。

\item \textbf{收缩集中在前 24 h，随后近似定半径干燥。} 附件 2 显示 $R$ 在前 6 h 掉到 1.374 cm、前 24 h 到 1.204 cm、之后稳定在 1.198--1.200 cm（图 16a）。因此对流项（浓缩效应）只在 $t\lesssim24$ h 活跃，后半程（24--52 h）本质上是``半径 1.2 cm 的定域干燥''，此时问题 4 的模型才与问题 1/2 的固定域格式同构。

\item \textbf{判据仍由中心控制。} 表面在 12 h 已达 0.2145、18 h 达 0.1016（早已``干透''），但中心到 52.39 h 才降到 0.15。与问题 3 一样，``各处低于 0.15''等价于``中心低于 0.15''（V5）。

\item \textbf{对流项是浓缩而非扩散。} 图 17b 的归一化剖面显示，固定 $E$ 处的浓度在收缩期先被``抬升''（水分被压入更小的体积），随后才随干燥单调下降；这在物理坐标图 17a 中表现为剖面右端随 $R(t)$ 同步内移。这正是 $+(E R'/R)\partial C/\partial E$ 项的物理后果。
\end{enumerate}

\section{验证证据（V1--V7）}

全部原始输出见 \texttt{results/problem4\_verification.txt}（\texttt{python p4\_verify.py} 生成）。

\subsection{空间收敛（V1）}

\begin{center}
\begin{tabular}{cccc}
\toprule
$N$ & $t_\text{dry}$ / s & $t_\text{dry}$ / h & 相邻差 / min \\
\midrule
100 & 188\,520 & 52.3667 & — \\
200 & 188\,580 & 52.3833 & +1.00 \\
400 & 188\,640 & 52.4000 & +1.00 \\
800 & 188\,640 & 52.4000 & +0.00 \\
\bottomrule
\end{tabular}
\end{center}

$N=400$ 与 800 结果一致到秒级，\textbf{空间已收敛}。交付取 $N=800$。

\subsection{时间收敛（V2）}

\begin{center}
\begin{tabular}{ccc}
\toprule
$\Delta t$ / s & $t_\text{dry}$ / s & $t_\text{dry}$ / h \\
\midrule
120 & 188\,640 & 52.4000 \\
60 & 188\,640 & 52.4000 \\
30 & 188\,610 & 52.3917 \\
\bottomrule
\end{tabular}
\end{center}

$\Delta t = 60$ s 与 120 s 一致，$\Delta t = 30$ s 差 30 s（0.5 min，在时间离散误差内）。\textbf{时间已收敛}。

\subsection{对流项符号的物理检验（V3）}

\begin{center}
\begin{tabular}{ccc}
\toprule
对流项符号 & $t_\text{dry}$ / h & 物理解释 \\
\midrule
\textbf{+1（正确）} & \textbf{52.400} & 收缩向内输送水分 $\to$ 浓缩 $\to$ 减缓干燥 \\
0（无对流） & 50.833 & 仅 $1/R^2$ 缩放 \\
$-1$（错误） & 48.650 & 收缩向外输送 $\to$ 非物理地加速干燥 \\
\bottomrule
\end{tabular}
\end{center}

\textbf{$+$ 号（链式法则严格导出）使 $t_\text{dry}$ 增大 1.57 h，符合物理}；$-$ 号使 $t_\text{dry}$ 减小，物理错误。对流项占 $t_\text{dry}$ 约 3\%，\textbf{不可忽略}。这既验证了符号，也说明题目给出的 $-E R'/R$ 是笔误。

\subsection{收缩效应（V4）}

\begin{center}
\begin{tabular}{lc}
\toprule
情形 & $t_\text{dry}$ / h \\
\midrule
有收缩（附件 2，$R$ 2$\to$1.2 cm） & \textbf{52.400} \\
无收缩（$R$ 冻结 2 cm） & 129.083 \\
\bottomrule
\end{tabular}
\end{center}

$R^2$ 缩小 2.79 倍、$1/R^2$ 放大，$t_\text{dry}$ 缩短 2.46 倍，收缩效应显著。

\subsection{判据等价性（V5）}

全部 3145 个输出时刻中，$\max_r C$ \textbf{真实偏离中心（$>10^{-9}$）的时刻为 0}。另有 40 个时刻的 $\arg\max$ 落在非中心节点，经核查全部发生在 $t\lesssim10$ min 的初始段，此时 $C$ 在径向上仍为平台值 2.55，中心与内部节点之差仅 $10^{-15}$ 量级（浮点并列噪声）。\textbf{决定终止时间的是``最迟达标节点''，即中心}。

\subsection{退化一致性（V6）}

$R'=0$ 时对流算子 $|\cdot|_{\max}=0$（精确为 0），Landau 格式\textbf{逐位退化}为固定域格式，不改变问题 1/2 的任何结果。

\subsection{离散质量守恒恒等式（V7）}

固定域 $R'=0$ 下，守恒型有限体积 + 表面 Robin 拆分给出离散恒等式 $\dfrac{\dd M}{\dd t}=-\dfrac{h_m}{R}(C_N-C_\infty)$（$M=\sum_jV_jC_j$）。前 2 h 内数值与解析右端之差 $\max=1.212\times10^{-5}$，即\textbf{质量守恒由格式结构精确保证}（残差来自 Picard 迭代容差，非守恒缺陷）。

\section{假设、局限与尚未解决的问题}

\begin{enumerate}
\item \textbf{半径 $R(t)$ 作为已知输入，而非双向耦合。} 本文把附件 2 的 $R(t)$ 当作给定边界直接使用，没有建立``水分流失量 $\to$ 收缩量''的本构关系。附件 2 是实测/经验数据，故移动边界本身是可靠的；但若要在未给数据的情形下外推，需要补收缩本构。
\item \textbf{移动边界模型仍是一维轴对称。} 收缩只体现在径向 $R(t)$，轴向长度变化与端部效应未建模（题目未给长度，沿用问题 1/2 口径）。
\item \textbf{$R'(t)$ 的端点单侧差分。} 附件 2 数据端点（$t=0$ 与 $t=259\,200$ s）用单侧差分，中段用中心差分。$R'$ 的初值（收缩最陡段）恰在端点附近，单侧差分的精度略低于中心差分；好在收缩在 $t=0$ 附近变化平滑，且对 $t_\text{dry}$ 的影响在分钟以下量级（V1/V2 已覆盖）。
\item \textbf{对流项与扩散项的同点离散。} 对流项中心差分在收缩最陡的前几分钟（$|v_E|$ 最大）可能引入轻微的非单调性，但 V5 显示它没有造成 $\max_r C$ 的真实偏离，且对 $t_\text{dry}$ 无影响。
\item \textbf{判据取严格不等式 + 分钟向上取整。} 与问题 3 一致：连续口径 52.3917 h，报出保守的 3144 min = 52.40 h。
\end{enumerate}

\section{可复现性}

\begin{verbatim}
cd code
python p4_model.py    # 单跑求解器（N=200，约 10 s）
python p4_run.py      # 交付口径（N=800），生成 result4.xlsx 与表 6
python p4_verify.py   # V1--V7 验证套件，生成 problem4_verification.txt
python p4_figures.py  # 四张配图（依赖 p4_fields.npz）
\end{verbatim}

依赖：\texttt{numpy}、\texttt{openpyxl}、\texttt{matplotlib}（不含 \texttt{scipy} 依赖）。

本次交付的实际运行环境与耗时：交付口径求解约 1 min（$N = 800$、3144 步）；验证套件约 3 min（含 V4 的 129 h 无收缩长算例）；绘图约 15 s。

\end{document}
